%% ---------------------------------------------------------------------------
\section{Discussion}\label{sec:discussion}

\paragraph{Answer to the research question} Within the design space it was trained on, the surrogate predicts member forces accurately enough to replace most repeated linear analyses at the level of preliminary design and screening: the median error of the end-moment field is 2.4\%, the largest end moment of a frame is within 0.8\% at the median, and the forces are exactly in equilibrium with the loads. Displacements are less accurate (median 4.5\% for the horizontal displacement field, 4.9\% for the governing drift). Outside the training space the median errors roughly triple and the 95th percentile of the moment error exceeds 20\%; the surrogate is then usable only together with a mechanism that sends doubtful cases back to the FE solver. That mechanism works for unseen geometries, where the ensemble spread ranks the errors well, but not for unseen load patterns, where the ensemble is confidently wrong. The answer is therefore a qualified yes: the surrogate can reduce repeated FE analysis inside a validated domain and, with referral, for moderately larger frames, but it is not a substitute for FE analysis on load patterns that the training data did not cover.

\paragraph{What the constraints buy} The constraints that hold exactly by construction contribute in two different ways. Exact load-case superposition and dimensional scaling are an \emph{accuracy} device: they remove the need to learn the dependence on the load intensities and on $E$, and they make the lateral and gravity responses equally well conditioned, which roughly halves the displacement errors (Table~\ref{tab:ablation}). The equilibrium projection is primarily an \emph{admissibility} device: it guarantees that the forces are in equilibrium, which a soft penalty does not achieve, and with a variance-weighted metric it also improves the forces, strongly for the axial forces and modestly for the moments. With a poorly chosen metric the projection can degrade the moments, so the metric is an essential part of the method rather than a detail. The projection does not enforce compatibility or the constitutive law: the predicted displacements and forces are not guaranteed to be consistent with each other, and imposing both equilibrium and compatibility exactly would amount to solving the FE problem.

\paragraph{Uncertainty as a decision tool} The two uses of uncertainty behaved differently. At deployment, the ensemble spread is a good indicator of which frames to refer to FE analysis for in-distribution and geometrically extrapolated frames, but its calibration collapses under a change of load pattern (coverage of the $\pm2\sigma$ intervals 67\% instead of 92\%). During training-set construction, greedy acquisition by ensemble variance was worse than random sampling, whereas acquisition by the equilibrium residual, a physics-based indicator that requires no labels, was the best rule for the member forces and halved the FE budget needed to reach a given accuracy on unseen geometries. A plausible explanation is that the residual is driven by the parts of the response that statics can check, whereas the variance also rewards the most atypical parameter combinations, which contribute little to the bulk of the test distribution. Batch-diversity and density-weighted acquisition rules, which are common remedies, were not tested.

\paragraph{Cost} The per-frame speed-up of the batched surrogate over PyNite is about 10--15, and 3.3--3.8 with 20\% referral. Because a linear PyNite analysis of these frames takes only 30--100\,ms, the one-off training cost of the ensemble (about 41\,min of CPU time) is recovered only after tens of thousands of evaluations. This is an honest limitation of applying the method to cheap linear analyses and should be weighed against the specific use case; the accuracy, admissibility and uncertainty results do not depend on it.

\paragraph{Limitations} The study is restricted to planar, regular, multi-storey moment frames with prismatic members, rigid joints, uniformly distributed beam loads and floor-level lateral forces, analysed with first-order linear elasticity. Sections and loads were sampled independently, so many frames would not satisfy serviceability limits; the dimensionless formulation makes the relative errors invariant to this, but the parameter correlations of designed frames are not represented. All results come from one dataset, one training run of each main model and two seeds of the adaptive study, so small differences (for example between the two projection metrics on in-distribution frames) should not be over-interpreted. The adaptive study is retrospective: candidates were labelled in advance and their labels revealed only on request, which yields exact budget accounting but not the scheduling overheads of a live FE loop. Timings were measured on one CPU with Python implementations of both the surrogate and PyNite.

\paragraph{Outlook} Natural extensions are a mixed formulation that also enforces compatibility through the member flexibilities, extension to three-dimensional frames, braced systems and second-order (P-$\Delta$) analysis, in which the FE cost is higher and superposition no longer holds exactly, explicit domain checks on the load pattern to complement the ensemble, and diversity-aware batch acquisition.

%% ---------------------------------------------------------------------------
\section{Conclusions}\label{sec:conclusions}

A mechanics-constrained surrogate modelling framework for the linear response of planar moment frames was developed and evaluated on 7{,}200 PyNite analyses. A graph network ensemble predicts nodal displacements and member basic forces in a dimensionless, load-case-separated form, so that superposition, homogeneity in the loads and in Young's modulus and the boundary conditions hold exactly; an equilibrium projection with an uncertainty-weighted metric makes the predicted forces statically admissible; and the uncertainty estimate decides which frames are analysed with PyNite during training and at deployment. The main conclusions are:
\begin{enumerate}
\item On in-distribution frames the surrogate reaches median errors of 2.4\% for the end-moment field, 0.6\% for the axial forces, 0.8\% for the largest end moment and 4.5\% for the horizontal displacements, four to six times lower than the best of five regression and ML baselines given identical features.
\item On unseen geometries and unseen load patterns the median moment errors are 7.5\% and 8.7\% (best baseline 22.1--22.5\% and 16.8--17.3\%), but the 95th percentiles exceed 20\%.
\item The equilibrium projection reduces the out-of-balance force from 4--16\% (raw ensemble) and up to 120\% (baselines) to machine precision for every model; with a variance-weighted metric it also reduces the force errors, whereas a Euclidean metric can increase the moment errors. A soft equilibrium penalty is ineffective.
\item Exact load-case superposition roughly halves the displacement errors compared with a single load scale.
\item The ensemble spread ranks the frame errors well for in-distribution frames and unseen geometries (Spearman 0.81 and 0.92) and supports effective FE referral there, but it is poorly calibrated for unseen load patterns.
\item For adaptive data generation, acquisition by the equilibrium residual of the raw prediction gave the most accurate member forces and reached the accuracy of 1{,}200 random analyses on unseen geometries with 600 analyses, whereas acquisition by ensemble variance was worse than random sampling.
\item The batched surrogate is about 10--15 times faster than a PyNite analysis per frame, but the training cost is recovered only after tens of thousands of evaluations for analyses as cheap as those considered here.
\end{enumerate}
The framework is therefore best viewed as a reliability layer that makes learned structural surrogates statically admissible and self-aware of their limits, rather than as a replacement for FE analysis.
